\chapter{documentatie vormen}
\label{ch:documentatie_vormen}

Om de tekortkomingen van traditionele documentatie te overbruggen, wordt er in dit onderzoek gekeken naar alternatieve vormen die de scheiding tussen uitleg en uitvoering opheffen. De meest veelbelovende benaderingen zijn geworteld in de filosofie van \textit{Literate Programming} en de moderne praktische uitwerking daarvan in de vorm van computationele notebooks. Deze vormen transformeren documentatie van een statisch bijproduct naar een actief, uitvoerbaar en centraal instrument binnen de wetenschappelijke workflow.

\section{Het fundament: Literate Programming}
Zoals geïntroduceerd door \textcite{Knuth1984}, draait Literate Programming om de premisse dat een computerprogramma in de eerste plaats geschreven moet worden als een narratief for menselijke lezers, waarbij de machine-uitvoerbare code een secundair derivaat is. In plaats van tekstuele uitleg als secundaire commentaar toe te voegen aan de code (\textit{code with comments}), vormt bij Literate Programming de tekst het primaire raamwerk waarin de codefragmenten zijn ingebed.

De essentie van Literate Programming ligt in de twee fundamentele operaties: \textbf{tangle} en \textbf{weave}. 
\begin{itemize}
    \item \textbf{Tangle:} Deze operatie extraheert de broncode uit het bronbestand en ordent deze in een formaat dat door de compiler of interpreter kan worden uitgevoerd. Voor EMAV betekent dit dat de wetenschappelijke formules die in een narratief document zijn beschreven, direct worden omgezet naar de C\#-logica die de emissieberekeningen aandrijft. Dit garandeert dat de "as-documented" logica identiek is aan de "as-implemented" logica.
    \item \textbf{Weave:} Deze operatie genereert de leesbare documentatie (zoals een PDF of HTML-pagina) uit hetzelfde bronbestand. Hierbij worden de codefragmenten gepresenteerd in hun context, vaak vergezeld van rijke visualisaties, wiskundige notaties in \LaTeX \@ en uitgebreide kruisverwijzingen.
\end{itemize}
Voor het EMAV-project biedt dit perspectief een structurele oplossing voor de \textit{documentation drift}. Wanneer de wetenschappelijke verantwoording en de eigenlijke berekeningslogica in hetzelfde artefact zijn ondergebracht, verdwijnt de noodzaak voor handmatige synchronisatie. Het document fungeert als een "Executable Paper", waarbij elke bewering over het model direct kan worden gestaafd door de bijbehorende code uit te voeren.

\section{Vergelijking van Computationele Notebook-platforms}
Computationele notebooks zijn de moderne, interactieve erfgenamen van Knuths visie. Ze bieden een "REPL-achtige" (\textit{Read-Eval-Print Loop}) ervaring die ideaal is voor exploratief onderzoek. Voor de implementatie binnen ILVO zijn drie belangrijke platforms overwogen, elk met hun eigen sterktes en zwaktes binnen een .NET-georiënteerde omgeving.

\subsection{Jupyter Notebooks}
Jupyter \autocite{jupyternotebook} is de \textit{de facto} standaard in de data science wereld. Het blinkt uit door zijn enorme ecosysteem en de ondersteuning voor honderden "kernels". In dit project is gekozen voor Jupyter vanwege de volwassenheid van de webinterface en de uitstekende integratie met Docker-omgevingen. 
Een specifiek voordeel voor EMAV is de mogelijkheid om via de \textbf{Python.NET} bridge de C\#-kern direct aan te spreken vanuit een Python-kernel. Dit laat onderzoekers toe om Python's krachtige visualisatie-libraries (zoals Matplotlib, Seaborn of Plotly) te gebruiken op data die is gegenereerd door de robuuste .NET-berekeningsengine. Het grootste nadeel van Jupyter is de niet-lineaire uitvoervolgorde van cellen, wat discipline vereist van de onderzoeker om de reproduceerbaarheid te waarborgen \autocite{Rule2018}.

\subsection{Quarto}
Quarto \autocite{quarto2025} is een technisch meer geavanceerde opvolger die specifiek focust op wetenschappelijke publicaties en technische rapportage. Het grote voordeel van Quarto is dat het gebruikmaakt van platte tekstbestanden (Markdown), wat het versiebeheer via Git aanzienlijk vergemakkelijkt in vergelijking met de JSON-gebaseerde \texttt{.ipynb} bestanden van Jupyter. Quarto ondersteunt ook "multi-language" documenten, waarbij verschillende talen in één flow kunnen worden gecombineerd. Hoewel Quarto een sterke kandidaat is voor de finale publicatie van wetenschappelijke rapporten binnen ILVO, is de interactieve "cell-by-cell" ervaring voor de gemiddelde onderzoeker momenteel nog minder intuïtief dan bij de klassieke Jupyter interface.

\subsection{Org-mode Babel}
Org-mode \autocite{orgmode2025} biedt wellicht de meest pure vorm van Literate Programming binnen de Emacs-omgeving. Met de Babel-extensie kunnen codeblokken in nagenoeg elke taal worden uitgevoerd en kunnen resultaten direct in het document worden teruggekoppeld. Voor de doeleinden van EMAV is Org-mode echter minder geschikt als breed verspreide oplossing. De steile leercurve van Emacs vormt een te grote barrière voor de gemiddelde wetenschapper die snel een berekening wil valideren. Het blijft echter een krachtige tool voor "power users" en ontwikkelaars die een maximale controle over hun documentatieproces wensen.

\section{De Techniek achter de Verweving: C\# Interop in Notebooks}
Een cruciaal aspect van de gekozen documentatievorm is de manier waarop de complexe .NET-codebase toegankelijk wordt gemaakt in de interactieve notebook-omgeving. In plaats van de logica te herimplementeren in een scripttaal (wat de betrouwbaarheid zou ondermijnen), wordt de eigenlijke \texttt{ILVO.EMAV.Core.dll} dynamisch ingeladen.

Dit proces verloopt via een startup-script dat gebruikmaakt van de \textbf{Common Language Runtime (CLR)}. Hierdoor kunnen onderzoekers in hun notebook direct objecten instantiëren van C\#-klassen, methoden aanroepen en zelfs LINQ-queries uitvoeren op de data-modellen. Deze aanpak garandeert dat de documentatie altijd werkt op de \textit{exacte} versie van de code die ook in de productieomgeving van EMAV-web draait. Het opheffen van de technische barrière tussen de gecompileerde C\#-wereld en de interactieve data science-wereld is de belangrijkste stap in het verkleinen van de kloof tussen IT en wetenschap.

\section{Interactieve Documentatie als Leer- en Validatieplatform}
De integratie van notebooks transformeert de documentatie van een passief naslagwerk naar een actieve leeromgeving. Naast de pure berekeningslogica kunnen notebooks worden verrijkt met interactieve elementen zoals sliders voor parameter-exploratie, dynamische grafieken en directe links naar externe wetenschappelijke bronnen. 

Dit zorgt voor een gelaagde informatievoorziening die aansluit bij verschillende behoeften:
\begin{enumerate}
    \item \textbf{De Narratieve Laag:} De wetenschappelijke tekst die de context, de wiskundige afleidingen en de verantwoording van het model biedt.
    \item \textbf{De Executabele Laag:} De levende codefragmenten die de formules direct implementeren en uitvoerbaar maken.
    \item \textbf{De Visualisatielaag:} Grafieken en tabellen die de output van de code direct inzichtelijk maken en vergelijking met historische data mogelijk maken.
\end{enumerate}
Door deze drie lagen in één document te verenigen, wordt voldaan aan de behoefte van de onderzoeker aan zowel transparantie als autonomie. Het resultaat is een documentatievorm die niet alleen beschrijft hoe het systeem werkt, maar die zelf een integraal en onmisbaar onderdeel vormt van het kwaliteitsborgingsproces van het instituut.
